% !TEX program = xelatex
\documentclass[11pt,a4paper]{ctexart}

\usepackage{amsmath,amssymb,amsfonts}
\usepackage{booktabs}
\usepackage{geometry}
\usepackage{hyperref}
\usepackage{enumitem}
\usepackage{xcolor}
\usepackage{longtable}

\geometry{left=2.5cm,right=2.5cm,top=2.5cm,bottom=2.5cm}
\hypersetup{colorlinks=true,linkcolor=blue,urlcolor=blue}

\title{%
  \textbf{星闪 SLB 物理层}\\[0.4em]
  \Large 6.2.1 一般描述技术文档\\[0.3em]
  \large 120\,kHz 基础子载波间隔系统 \textbullet\ 物理层内容分类与信息块
}
\author{依据 T/XS 10001-2025《星闪无线通信系统 接入层\\
同步低时延宽带空口 SLB 技术要求和测试方法》整理\\
（团体标准 20250326 版）}
\date{2026 年 6 月 22 日}

\begin{document}
\maketitle
\tableofcontents
\newpage

%======================================================================
\part{总览}
%======================================================================

\section{文档范围与模块划分}

本文档对应 SLB 120\,kHz 物理层协议第 6.2.1 节``一般描述''，为后续 6.2.2\textasciitilde 6.2.10 各子节的工程实现与联调提供\textbf{统一的物理层内容分类索引}与\textbf{通信域/直通链路 PIB 划分}参考。

6.2.1 本身不规定波形参数、帧结构或编码算法，而是定义 120\,kHz 子载波间隔系统中物理层传输内容的\textbf{三类集合}（信号、控制信息、数据信息）及其在\textbf{物理层信息块（PIB）}中的打包方式。

\begin{table}[h]
\centering
\small
\begin{tabular}{lll}
\toprule
\textbf{分类} & \textbf{数量} & \textbf{后续详细条款} \\
\midrule
物理层信号 & 8 类 & 6.2.3 物理层信号 \\
物理层控制信息 & 17 类 & 6.2.4 物理层控制信息 \\
物理层数据信息 & 2 类 & 6.2.5 物理层数据信息 \\
物理层信息块（PIB） & 通信域 4 块 / 直通 2 块 & 6.2.3\textasciitilde 6.2.5 各块内映射 \\
\bottomrule
\end{tabular}
\end{table}

\section{与 6.1 物理层一般描述的关系}

6.1 节从系统层面定义物理层向数据链路层提供的服务（CRC/FEC/HARQ、速率匹配、资源映射、调制解调、同步、测量、MIMO、波束赋形、射频处理等），并给出 CP-OFDM+TDD 的总体框架及三类内容的通用命名规则（G/T/附/直通链路前缀约定）。

6.2.1 在 6.1 框架下，针对\textbf{120\,kHz 基础子载波间隔系统}，枚举该 SCS 下具体包含的信号、控制、数据类型及 PIB 组成，是 6.2 节的入口索引。

\section{120\,kHz 子载波间隔系统定位}

\begin{table}[h]
\centering
\begin{tabular}{ll}
\toprule
\textbf{概念} & \textbf{说明} \\
\midrule
子载波间隔 120\,kHz & $\Delta f = f_s/256 = 120$\,kHz，$f_s = 30.72$\,MHz \\
资源元素（RE） & 时频网格最小映射单位，详见 6.2.2 \\
物理层三类内容 & 信号（辅助解调/同步/测量）、控制信息（调度/接入/反馈）、数据信息（用户载荷） \\
\bottomrule
\end{tabular}
\end{table}

\section{后续章节索引}

\begin{table}[h]
\centering
\small
\begin{tabular}{ll}
\toprule
\textbf{节号} & \textbf{内容} \\
\midrule
6.2.2 & 基础帧结构和物理资源（CP-OFDM 波形、超帧/无线帧、RE 网格） \\
6.2.3 & 物理层信号（同步、DMRS、CSI-RS、PMS 等生成与映射） \\
6.2.4 & 物理层控制信息（SAB/RAB/G-PCIB 等块内字段） \\
6.2.5 & 物理层数据信息（第一/二类数据 TB 处理） \\
6.2.6 & 信道编码（极化码/LDPC） \\
6.2.7\textasciitilde 6.2.10 & 同步、基带生成、调制上变频 \\
\bottomrule
\end{tabular}
\end{table}

\newpage
%======================================================================
\part{6.2.1 一般描述}
%======================================================================

\section{协议章节对应}

本节对应星闪 SLB 接入层物理层协议 T/XS 10001-2025 第 6.2 节``120\,kHz 基础子载波间隔系统''之 \textbf{6.2.1 一般描述}，完整涵盖：
\begin{itemize}[nosep]
  \item 120\,kHz 系统中物理层信号清单（8 类）；
  \item 物理层控制信息清单（17 类）；
  \item 物理层数据信息清单（2 类）；
  \item 物理层信息块（PIB）定义；
  \item 通信域与直通链路的 PIB 组成及散布资源。
\end{itemize}

\section{功能定义}

6.2.1 的核心功能为：在 120\,kHz SCS 下建立物理层传输内容的\textbf{分类体系}与\textbf{信息块打包规则}，使接收端能够：
\begin{enumerate}[nosep]
  \item 识别当前 RE 上承载的是信号、控制还是数据；
  \item 按 PIB 类型（SAB/RAB/G-PCIB 等）确定块边界与解码流程；
  \item 区分通信域（经 G 节点）与直通链路（设备直连）的资源组织差异。
\end{enumerate}

\begin{table}[h]
\centering
\small
\begin{tabular}{p{3.5cm}p{4cm}p{6cm}}
\toprule
\textbf{内容类型} & \textbf{是否承载高层信息} & \textbf{典型作用} \\
\midrule
物理层信号 & 否 & 同步、信道估计参考、探测/CSI、相位辅助、定位测量 \\
物理层控制信息 & 是（控制面） & 接入、调度、HARQ/CQI 反馈、功耗与载波管理 \\
物理层数据信息 & 是（用户面） & 第一/二类业务载荷 \\
物理层信息块 & 组合上述三类 & 标准化时频打包，便于接收端块级处理 \\
\bottomrule
\end{tabular}
\end{table}

\section{模块边界（上下游及职责划分）}

\subsection{上游模块（输入来源）}

\begin{table}[h]
\centering
\small
\begin{tabular}{p{4.2cm}p{4.5cm}p{5.5cm}}
\toprule
\textbf{上游模块} & \textbf{输入给本模块的内容} & \textbf{说明} \\
\midrule
数据链路层（MAC/RLC） & 待发送的控制/数据比特 & 由 6.2.4/6.2.5 组装为物理层信息 \\
帧配置/调度模块 & 当前域类型（通信域/直通）、链路类型 & 决定 PIB 选择与 RE 映射 \\
6.2.2 帧结构模块 & 时频资源网格 & RE 位置由超帧/无线帧/符号配置决定 \\
\bottomrule
\end{tabular}
\end{table}

\subsection{下游模块（输出去向）}

\begin{table}[h]
\centering
\small
\begin{tabular}{p{4.2cm}p{4.5cm}p{5.5cm}}
\toprule
\textbf{下游模块} & \textbf{本模块输出} & \textbf{说明} \\
\midrule
6.2.3 物理层信号 & 信号类型标识 + RE 映射需求 & 同步/DMRS/CSI-RS/PMS 等 \\
6.2.4 物理层控制信息 & 控制信息类型 + PIB 归属 & SAB/RAB/G-PCIB/P-PCIB 等 \\
6.2.5 物理层数据信息 & 数据类型（第一/二类）+ RE 映射 & 与 DMRS-D 绑定 \\
6.2.6/6.5 编码与通用功能 & 待编码比特流 & CRC、极化码、加扰、调制 \\
\bottomrule
\end{tabular}
\end{table}

\subsection{本模块不负责的内容}

\begin{table}[h]
\centering
\small
\begin{tabular}{p{5.5cm}p{8.5cm}}
\toprule
\textbf{不负责内容} & \textbf{原因} \\
\midrule
CP-OFDM 波形参数与帧结构 & 由 6.2.2 定义 \\
各信号/控制/数据的具体比特字段 & 由 6.2.3\textasciitilde 6.2.5 各子节定义 \\
信道编码与调制算法 & 由 6.2.6、6.5 节定义 \\
射频发射与功控 & 由 6.2.10、6.5.4 及射频章定义 \\
\bottomrule
\end{tabular}
\end{table}

\section{在 SLB 通信流程中的角色}

6.2.1 处于 PHY 层内容组织的\textbf{最顶层分类}，在发送/接收流水线中位于帧结构（6.2.2）与具体信号/控制/数据生成（6.2.3\textasciitilde 6.2.5）之间，提供``传什么类型内容、打包成哪个 PIB''的语义框架。

\textbf{通信域终端接收故事线：}
\begin{enumerate}[nosep]
  \item 搜同步信号 $\rightarrow$ 读 SAB 中的同步信息（定时、帧、标识）；
  \item 读广播 $\rightarrow$ 获知组/网络能力与部分配置；
  \item 要接入 $\rightarrow$ 看 RAB 的接入信息，发起接入；
  \item 接入后 $\rightarrow$ 解析 G-PCIB（及 G-DMRS-C、T 控制 DMRS）获取调度、HARQ、CQI、SR、休眠唤醒、载波切换等；
  \item 传数据 $\rightarrow$ 按调度映射第一/二类数据 + DMRS-D + 必要时相位调整；附链路数据可在 A-PDIB；
  \item 链路优化 $\rightarrow$ 信道探测/CSI 参考 $\rightarrow$ 终端上报 CQI；调度器调整 MCS；
  \item 定位/感知 $\rightarrow$ 按 G/T 位置测量资源指示使用位置测量信号；
  \item 未调度 RE $\rightarrow$ 空闲，不发送或仅接收噪声。
\end{enumerate}

\textbf{直通链路简化流程：}SAB $\rightarrow$ P-PCIB 调度 $\rightarrow$ P-DMRS-C + DMRS-D + 数据。

\section{强制要求与关键参数（T/XS 10001-2025 原文）}

以下内容为标准 6.2.1 节\textbf{完整原文}，与团体标准 20250326 版一致。

\subsection{6.2.1 一般描述}

120\,kHz 基础子载波间隔系统中，物理层信号包括：
\begin{itemize}[nosep]
  \item 同步信号
  \item G 链路公共解调参考信号
  \item T 链路控制信息解调参考信号
  \item 数据信息解调参考信号
  \item 数据信息相位调整信号
  \item 信道探测信号
  \item 信道状态信息参考信号
  \item 位置测量信号
\end{itemize}

物理层控制信息包括：
\begin{itemize}[nosep]
  \item 同步信息
  \item 广播信息
  \item 接入信息
  \item 动态调度数据控制信息
  \item 预配置调度激活去激活信息
  \item 休眠与唤醒指示信息
  \item 快速载波切换指示信息
  \item 初始接入载波指示信息
  \item G 链路位置测量信号或感知测量信号资源指示信息
  \item T 链路位置测量信号资源指示信息
  \item 直通链路动态调度数据控制信息
  \item 单级 HARQ 反馈信息
  \item 两级 HARQ 反馈信息
  \item CQI 反馈信息
  \item 调度请求信息
\end{itemize}

物理层数据信息包括：
\begin{itemize}[nosep]
  \item 第一类物理层数据信息
  \item 第二类物理层数据信息
\end{itemize}

用于实现特定功能，映射多种物理层信号、物理层控制信息、物理层数据信息组合的一系列资源元素的集合称为\textbf{物理层信息块}。

\textbf{通信域}物理层信息块包括：同步信息块 SAB、接入信息块 RAB、G 链路物理层控制信息块 G-PCIB、附链路物理层数据信息块 A-PDIB。

此外，通信域上还包括映射物理层广播信息、G 链路公共解调参考信号 G-DMRS-C、T 链路控制信息、数据信息解调参考信号 DMRS-D、物理层数据信息、信道探测信号、信道状态信息参考信号、位置测量信号和物理层数据信息相位调整信号的资源，以及未分配或未调度的空闲资源。

\textbf{直通链路}物理层信息块包括：同步信息块 SAB、直通链路物理层控制信息块 P-PCIB。

此外，直通链路上还包括映射直通链路公共解调参考信号 P-DMRS-C、数据信息解调参考信号 DMRS-D、物理层数据信息的资源，以及未分配或未调度的空闲资源。

\section{物理层信号详解（8 类）}

\begin{table}[h]
\centering
\small
\begin{tabular}{p{1cm}p{4.5cm}p{8cm}}
\toprule
\textbf{序号} & \textbf{信号名称} & \textbf{功能说明} \\
\midrule
1 & 同步信号 & 时间/频率对齐，发现网络，读取同步信息块 SAB \\
2 & G 链路公共解调参考信号（G-DMRS-C） & G 链路组/网关侧公共参考；接收端估计信道后解调控制/数据 \\
3 & T 链路控制信息解调参考信号 & 为 T 链路上控制信息提供 DMRS，与控制 RE 绑定 \\
4 & 数据信息解调参考信号（DMRS-D） & 为第一/二类物理层数据提供参考，支撑 MIMO、均衡与解调 \\
5 & 数据信息相位调整信号 & 补偿相位误差（频偏、预编码、相位跟踪等） \\
6 & 信道探测信号 & 主动或半主动探测信道时变、多径与干扰，为调度提供输入 \\
7 & 信道状态信息参考信号 & 供接收端测量信道质量并上报 CQI，调度器据此选择 MCS 与资源 \\
8 & 位置测量信号 & 支持定位/测距，与 G/T 链路位置测量资源指示控制信息配合 \\
\bottomrule
\end{tabular}
\end{table}

\textbf{小结}：信号 = 同步 + 参考（DMRS）+ 探测/CSI + 相位辅助 + 定位测量。详细生成规则见 6.2.3。

\section{物理层控制信息详解（17 类）}

\subsection{网络发现与接入}

\begin{table}[h]
\centering
\small
\begin{tabular}{ll}
\toprule
\textbf{控制信息} & \textbf{作用} \\
\midrule
同步信息 & 帧结构、组/小区标识、定时关系等，多在 SAB 中 \\
广播信息 & 全组可见的参数与能力概要 \\
接入信息 & 随机/竞争接入参数，对应 RAB \\
初始接入载波指示信息 & 指示终端首先在哪个载波上接入 \\
\bottomrule
\end{tabular}
\end{table}

\subsection{调度与资源}

\begin{table}[h]
\centering
\small
\begin{tabular}{ll}
\toprule
\textbf{控制信息} & \textbf{作用} \\
\midrule
动态调度数据控制信息 & 实时指示 UE 在何时何频传何种数据（通信域） \\
预配置调度激活去激活信息 & 半静态 grant，用短信令开/关，降低时延 \\
直通链路动态调度数据控制信息 & 直通链路上的同类调度 \\
G 链路位置测量信号或感知测量信号资源指示信息 & 告知 G 侧测量或感知 RE 位置 \\
T 链路位置测量信号资源指示信息 & 告知 T 侧定位资源位置 \\
\bottomrule
\end{tabular}
\end{table}

\subsection{功耗与载波}

\begin{table}[h]
\centering
\small
\begin{tabular}{ll}
\toprule
\textbf{控制信息} & \textbf{作用} \\
\midrule
休眠与唤醒指示信息 & 低功耗：何时进入休眠、何时唤醒 \\
快速载波切换指示信息 & 多载波场景下快速换频，支撑宽带与抗干扰 \\
\bottomrule
\end{tabular}
\end{table}

\subsection{反馈与请求（闭环链路自适应）}

\begin{table}[h]
\centering
\small
\begin{tabular}{ll}
\toprule
\textbf{控制信息} & \textbf{作用} \\
\midrule
单级 HARQ 反馈信息 & 一次 ACK/NACK，结构简单 \\
两级 HARQ 反馈信息 & 两阶段反馈，兼顾时延与可靠性 \\
CQI 反馈信息 & 信道质量指示，供调度选择速率 \\
调度请求信息 & 终端``我有数据要发''的上行请求 \\
\bottomrule
\end{tabular}
\end{table}

详细字段定义见 6.2.4 各子节。

\section{物理层数据信息详解（2 类）}

\begin{table}[h]
\centering
\small
\begin{tabular}{ll}
\toprule
\textbf{类型} & \textbf{说明} \\
\midrule
第一类物理层数据信息 & 主业务载荷的一种编码/映射类别 \\
第二类物理层数据信息 & 另一类载荷，可对应不同 QoS 或链路绑定 \\
\bottomrule
\end{tabular}
\end{table}

两类数据均依赖 DMRS-D 进行解调；具体调制编码、TB 分段见 6.2.5 与 6.2.6。

\section{物理层信息块（PIB）与两域对比}

\subsection{PIB 定义}

为实现特定功能，将多种物理层信号、控制信息、数据信息组合映射到一系列资源元素（RE）上的集合。PIB 是``标准化打包的时频包裹''，接收端先识别块类型与边界，再按块内规则解码。

\subsection{通信域 PIB}

\begin{table}[h]
\centering
\small
\begin{tabular}{lll}
\toprule
\textbf{PIB} & \textbf{全称} & \textbf{主要承载} \\
\midrule
SAB & 同步信息块 & 同步信号 + 同步信息 \\
RAB & 接入信息块 & 接入信息及相关内容 \\
G-PCIB & G 链路物理层控制信息块 & G 侧物理层控制（调度、指示、反馈等） \\
A-PDIB & 附链路物理层数据信息块 & 附链路上的数据信息 \\
\bottomrule
\end{tabular}
\end{table}

此外，通信域还包括单独映射（不整包进上述四块）的资源：广播信息、G-DMRS-C、T 链路控制信息、DMRS-D、物理层数据信息、信道探测信号、CSI 参考、位置测量信号、相位调整信号，以及未分配/未调度的空闲 RE。

\subsection{直通链路 PIB}

\begin{table}[h]
\centering
\small
\begin{tabular}{ll}
\toprule
\textbf{PIB} & \textbf{作用} \\
\midrule
SAB & 直通链路同步 \\
P-PCIB & 直通链路物理层控制（含直通动态调度等） \\
\bottomrule
\end{tabular}
\end{table}

此外还包括：P-DMRS-C（直通公共 DMRS）、DMRS-D、物理层数据信息、空闲 RE。

\subsection{通信域 vs 直通链路}

\begin{table}[h]
\centering
\small
\begin{tabular}{llll}
\toprule
\textbf{域} & \textbf{典型拓扑} & \textbf{主要 PIB} & \textbf{特点} \\
\midrule
通信域 & 设备经 G 节点接入，含 G/T/附链路 & SAB、RAB、G-PCIB、A-PDIB & 控制面完整：广播、接入、HARQ、CQI 等 \\
直通链路 & 设备之间直接通信 & SAB、P-PCIB & 结构精简，侧重直通调度与数据 \\
\bottomrule
\end{tabular}
\end{table}

\textbf{要点}：直通链路没有 RAB、G-PCIB、A-PDIB 等完整通信域结构，体现 P2P 精简、低时延设计。

\section{数据类型与内容维度}

\begin{table}[h]
\centering
\small
\begin{tabular}{llll}
\toprule
\textbf{内容类型} & \textbf{条目数} & \textbf{映射单位} & \textbf{是否进 PIB} \\
\midrule
物理层信号 & 8 & RE 集合 & 部分进 SAB 等，部分散布 \\
物理层控制信息 & 17 & RE 集合 & SAB/RAB/G-PCIB/P-PCIB 等 \\
物理层数据信息 & 2 & RE 集合 & G/T/附/直通链路 RE + A-PDIB \\
\bottomrule
\end{tabular}
\end{table}

\section{时序关系}

6.2.1 为静态分类定义，无时序状态机。各类内容在超帧/无线帧内的\textbf{出现顺序与周期}由 6.2.2 帧结构与 6.2.4 各控制信息调度规则共同决定。典型时序关系：
\begin{itemize}[nosep]
  \item SAB：每超帧/无线帧固定位置，终端首先接收；
  \item RAB：接入阶段按需；
  \item G-PCIB/P-PCIB：每 TTI 或预配置周期；
  \item 数据 + DMRS-D：按动态/预配置调度指示映射。
\end{itemize}

\section{链路调用对照表}

\begin{table}[h]
\centering
\footnotesize
\begin{tabular}{p{3.5cm}p{2.5cm}p{2.5cm}p{5cm}}
\toprule
\textbf{内容} & \textbf{通信域 PIB} & \textbf{直通链路 PIB} & \textbf{详细条款} \\
\midrule
同步信号 + 同步信息 & SAB & SAB & 6.2.3.1、6.2.4.1 \\
接入信息 & RAB & --- & 6.2.4.2 \\
G 侧控制（调度/HARQ/CQI 等） & G-PCIB & --- & 6.2.4.3\textasciitilde 6.2.4.10 \\
附链路数据 & A-PDIB & --- & 6.2.5 \\
直通调度 & --- & P-PCIB & 6.2.4 \\
G-DMRS-C / P-DMRS-C & 散布 & 散布 & 6.2.3.2 \\
DMRS-D & 散布 & 散布 & 6.2.3.3 \\
第一/二类数据 & 散布 & 散布 & 6.2.5 \\
\bottomrule
\end{tabular}
\end{table}

\section{与 5G/LTE 粗类比（实现参考，勿等同）}

\begin{table}[h]
\centering
\small
\begin{tabular}{ll}
\toprule
\textbf{星闪 120\,kHz 概念} & \textbf{粗类比} \\
\midrule
SAB & 类似 SSB / 同步块 \\
RAB & 类似 PRACH / 接入相关区域 \\
G-PCIB / P-PCIB & 类似 PDCCH / sidelink 控制区域 \\
G-DMRS-C / P-DMRS-C & 类似公共 DMRS \\
DMRS-D & 类似数据专用 DMRS \\
动态 + 预配置调度 & 类似 DCI + 配置 grant（减时延） \\
单级/两级 HARQ & HARQ 反馈粒度不同 \\
PIB & 星闪特有的时频打包结构 \\
\bottomrule
\end{tabular}
\end{table}

\section{约束与实现要点}

\begin{enumerate}[nosep]
  \item \textbf{域类型不可混用}：通信域 PIB 结构（RAB/G-PCIB/A-PDIB）不适用于直通链路；
  \item \textbf{链路前缀}：G/T/附/直通前缀按 6.1 命名规则使用，不产生混淆时可省略；
  \item \textbf{空闲 RE}：未分配/未调度 RE 不发送或仅接收噪声，接收端不得当作有效载荷解析；
  \item \textbf{散布资源}：广播、DMRS、数据等可不整包进 PIB，但须与 6.2.3\textasciitilde 6.2.5 映射规则一致；
  \item \textbf{比特顺序}：下标 0 为 LSB，空口先传低位（6.1 通用约定，后续各节均适用）。
\end{enumerate}

\section{速查表}

\subsection{物理层信号（8）}
同步 \textbar\ G-DMRS-C \textbar\ T 控制 DMRS \textbar\ DMRS-D \textbar\ 数据相位调整 \textbar\ 信道探测 \textbar\ CSI 参考 \textbar\ 位置测量

\subsection{物理层控制（17）}
同步 \textbar\ 广播 \textbar\ 接入 \textbar\ 动态调度 \textbar\ 预配置调度激活/去激活 \textbar\ 休眠唤醒 \textbar\ 快速载波切换 \textbar\ 初始接入载波 \textbar\ G 位置/感知资源指示 \textbar\ T 位置资源指示 \textbar\ 直通动态调度 \textbar\ 单级 HARQ \textbar\ 两级 HARQ \textbar\ CQI \textbar\ SR

\subsection{物理层数据（2）}
第一类 \textbar\ 第二类

\subsection{通信域 PIB}
SAB、RAB、G-PCIB、A-PDIB + 散布资源 + 空闲 RE

\subsection{直通链路 PIB}
SAB、P-PCIB + P-DMRS-C、DMRS-D、数据 + 空闲 RE

\section{与标准其他章节的衔接}

\begin{itemize}[nosep]
  \item \textbf{6.1}：物理层通用服务与 CP-OFDM 框架，6.2.1 为其在 120\,kHz SCS 下的内容枚举；
  \item \textbf{6.2.2}：RE 网格、超帧/无线帧、五种 CP-OFDM 符号类型，为 PIB 映射提供时频坐标；
  \item \textbf{6.2.3}：8 类信号的生成、序列与 RE 映射细节；
  \item \textbf{6.2.4}：17 类控制信息的编码、CRC、极化码与块内字段；
  \item \textbf{6.2.5}：第一/二类数据的 TB 处理、分段与映射；
  \item \textbf{6.2.6/6.5}：信道编码、CRC、加扰、调制、功控等通用算法；
  \item \textbf{6.3}：灵活带宽系统在相同内容分类下的带宽扩展；
  \item \textbf{第 8 章}：射频与信道参数，与载波切换、初始接入载波指示配合。
\end{itemize}

\vfill
\begin{center}
\small\textcolor{gray}{精确参数与字段定义以 T/XS 10001-2025 正式标准为准；本文档仅供 SLB 120\,kHz 物理层实现与联调参考。}
\end{center}

\end{document}
